\chapter{考点方向一：机器学习理论 (Machine Learning Theory)}
\label{ch:ml_theory}

\section{监督学习、PAC、VC 与泛化}

\teaching
监督学习的基本问题是：从有限样本中学一个预测规则\footnote{预测规则就是一个函数 \(h:\mathcal{X}\to\mathcal{Y}\)：看到新输入 \(x\) 时，输出 \(h(x)\)。形式上，先设一对随机变量 \((X,Y)\sim\Dcal\)，其中 \(X\) 是输入、\(Y\) 是答案；“训练集由 \(\Dcal\) 独立同分布抽样”是指每个 \((x_i,y_i)\) 都按同一个联合分布 \(P_{XY}\) 生成，且不同样本彼此独立。二分类中，\(x\) 可为邮件的词频和发件人特征，\(h(x)\in\{\text{垃圾},\text{正常}\}\)；回归中，\(x\) 可为房屋面积和位置，\(h(x)\in\R\) 预测房价；概率分类中，模型先输出 \(p_\theta(y\mid x)\)，再以 \(\argmax_y p_\theta(y\mid x)\) 作为类别预测。样本只是未知总体分布的一小批随机观测，因此训练误差小并不自动说明新样本上的误差也小。}，并希望它在未见样本上也表现好。这里的核心困难不是“训练集上能不能拟合”，而是“训练集误差能不能代表总体误差”。

设输入空间为 \(\mathcal{X}\)，标签空间为 \(\mathcal{Y}\)，未知数据分布为 \(\Dcal\)。训练集

\[
  S=((x_1,y_1),\ldots,(x_m,y_m))
\]

由 \(\Dcal\) 独立同分布抽样得到。假设类 \(\Hcal\) 是候选预测器集合，损失函数 \(\ell(h,(x,y))\) 衡量预测器 \(h\) 在样本 \((x,y)\) 上的错误。经验风险和真实风险分别为

\[
  L_S(h)=\frac1m\sum_{i=1}^{m}\ell(h,(x_i,y_i)),
  \qquad
  L_{\Dcal}(h)=\E_{(x,y)\sim\Dcal}\ell(h,(x,y)).
\]

经验风险是可计算的训练误差；真实风险是我们真正关心但无法直接计算的泛化误差。经验风险最小化（ERM）写作

\[
  h_S\in \argmin_{h\in\Hcal}L_S(h).
\]

这一定义背后的直觉很简单：既然看不到总体分布，就先选训练集上表现最好的模型。但 ERM 是否可靠取决于 \(\Hcal\) 的复杂度。如果 \(\Hcal\) 太大，它可以把随机噪声也拟合掉，训练误差会系统性偏低。

\subsection{PAC 学习}

\concepttarget{concept:pac}

PAC 是 Probably Approximately Correct 的缩写。它把“学会”写成一个概率命题：样本足够多时，学习算法以高概率输出一个误差不超过目标精度的假设。

在二分类可实现情形中，假设存在目标函数 \(f\in\Hcal\)，标签由 \(y=f(x)\) 生成。算法 \(A\) PAC 学会 \(\Hcal\)，意思是对任意精度 \(\varepsilon>0\) 和失败概率 \(\delta>0\)，当样本量

\[
  m\ge m_{\Hcal}(\varepsilon,\delta)
\]

时，有

\[
  \Pbb_{S\sim \Dcal^m}
  \left[
    L_{\Dcal}(A(S))\le \varepsilon
  \right]
  \ge 1-\delta.
\]

这里 \(\varepsilon\) 是 accuracy parameter，控制允许多错；\(\delta\) 是 confidence parameter，控制允许多大概率失败。样本复杂度 \(m_{\Hcal}(\varepsilon,\delta)\) 告诉我们要多少样本才能达到这个保证。

为什么 Bernoulli 和 union bound 会不断出现？对固定的 \(h\)，二分类 \(0\)-\(1\) 损失可写成误差指示变量

\[
  Z_i=\mathbb{I}\{h(x_i)\ne f(x_i)\}.
\]

于是

\[
  L_S(h)=\frac1m\sum_{i=1}^{m}Z_i,
  \qquad
  \E[Z_i]=L_{\Dcal}(h).
\]

固定 \(h\) 时，问题就是 Bernoulli 均值集中；若要同时控制所有 \(h\in\Hcal\)，就要把坏事件

\[
  \left|L_S(h)-L_{\Dcal}(h)\right|>\varepsilon
\]

对整个假设类做并集控制。这就是有限假设类界中 \(\log|\Hcal|\) 的来源。若 \(\Hcal\) 有限，Hoeffding 不等式和 union bound 给出典型形式

\[
  \Pbb\left[
    |L_S(h)-L_{\Dcal}(h)|>\varepsilon
  \right]
  \le
  2\exp(-2m\varepsilon^2)
\]

对每个固定 \(h\) 成立。因为

\[
  \left\{
    \exists h\in\Hcal:
    |L_S(h)-L_{\Dcal}(h)|>\varepsilon
  \right\}
  =
  \bigcup_{h\in\Hcal}
  \left\{
    |L_S(h)-L_{\Dcal}(h)|>\varepsilon
  \right\},
\]

所以由 union bound，

\[
  \Pbb\left[
    \exists h\in\Hcal:
    |L_S(h)-L_{\Dcal}(h)|>\varepsilon
  \right]
  \le
  2|\Hcal|\exp(-2m\varepsilon^2).
\]

令右边不超过 \(\delta\)，得到

\[
  m
  \ge
  \frac{1}{2\varepsilon^2}
  \log\frac{2|\Hcal|}{\delta}.
\]

这条公式不是为了背常数，而是为了看清结构：样本量随 \(1/\varepsilon^2\) 增长，随 \(\log(1/\delta)\) 增长，随模型类大小的对数增长。

\paragraph{可实现 ERM 的 PAC 推导}

考试中更常见的是可实现情形：算法输出任意训练误差为零的 \(h_S\)。要证明 \(h_S\) 真误差小，只需排除“坏但碰巧一致”的假设。称 \(h\) 是坏假设，如果

\[
  L_{\Dcal}(h)>\varepsilon.
\]

对固定坏假设 \(h\)，一个样本不暴露它错误的概率至多为 \(1-\varepsilon\)。由于样本独立，\(m\) 个样本都不暴露它错误的概率至多为

\[
  (1-\varepsilon)^m
  \le
  \exp(-m\varepsilon).
\]

若 \(\Hcal\) 有限，存在某个坏假设仍与训练集完全一致的概率不超过

\[
  |\Hcal|\exp(-m\varepsilon).
\]

令该概率不超过 \(\delta\)，得到可实现 PAC 的样本量条件

\[
  m
  \ge
  \frac{1}{\varepsilon}
  \log\frac{|\Hcal|}{\delta}.
\]

这条界和上面的二侧统一收敛界常数与 \(\varepsilon\) 阶不同，原因是它用了更强的“训练误差为零”和“可实现”假设 \parencite{shalev2014uml,mohri2018foundations}。

\examnote
考题常要求从“固定 \(h\)”过渡到“存在 \(h\)”：第一步写 Bernoulli 误差变量，第二步用集中不等式，第三步对 \(\Hcal\) 做 union bound。若漏掉第三步，通常只能得到单个假设的界。

\subsection{Agnostic PAC 与统一收敛}

\concepttarget{concept:agnostic-pac}

现实中目标函数未必属于 \(\Hcal\)，甚至标签可能含噪声。Agnostic PAC 不假设可实现性，而是要求算法输出接近类内最优的假设：

\[
  L_{\Dcal}(A(S))
  \le
  \inf_{h\in\Hcal}L_{\Dcal}(h)+\varepsilon
\]

以至少 \(1-\delta\) 的概率成立。这里 \(\inf_{h\in\Hcal}L_{\Dcal}(h)\) 是\emph{类内最优风险}（class-optimal risk），并不是 approximation error。相对于 Bayes 风险 \(L^*_{\mathrm{Bayes}}\) 的近似误差应写作 \(\inf_{h\in\Hcal}L_{\Dcal}(h)-L^*_{\mathrm{Bayes}}\)。下文若使用近似 ERM，还须把优化误差单列；本段先分析精确 ERM。ERM 的困难变成证明统一收敛：

\[
  \sup_{h\in\Hcal}
  |L_S(h)-L_{\Dcal}(h)|
  \le \varepsilon.
\]

一旦统一收敛成立，ERM 的泛化分析几乎是机械的。令 \(h^*\in\argmin_{h\in\Hcal}L_{\Dcal}(h)\)，则

\[
  L_{\Dcal}(h_S)
  \le L_S(h_S)+\varepsilon
  \le L_S(h^*)+\varepsilon
  \le L_{\Dcal}(h^*)+2\varepsilon.
\]

第一步和第三步来自统一收敛，第二步来自 ERM 定义。这类“三行证明”是学习理论真题最常见的骨架。

如果题目要求最终 excess risk 不超过 \(\varepsilon\)，通常把统一收敛精度设成 \(\varepsilon/2\)。即若

\[
  \sup_{h\in\Hcal}
  |L_S(h)-L_{\Dcal}(h)|
  \le \frac{\varepsilon}{2},
\]

则同样推导给出

\[
  L_{\Dcal}(h_S)
  \le
  L_{\Dcal}(h^*)+\varepsilon.
\]

这就是 agnostic ERM 样本复杂度里经常出现额外常数的来源，而不是新的理论困难。

\subsection{VC 维、增长函数与 Sauer 引理}

\concepttarget{concept:vc-sauer}



当 \(\Hcal\) 无限时，不能再用 \(|\Hcal|\) 直接度量复杂度。VC 理论度量的是假设类在有限点集上实现任意标记的能力。

\definitionlabel 给定 \(m\) 个点 \(x_1,\ldots,x_m\)，增长函数定义为

\[
  \tau_{\Hcal}(m)
  =
  \max_{x_1,\ldots,x_m}
  \left|
  \{
    (h(x_1),\ldots,h(x_m)):h\in\Hcal
  \}
  \right|.
\]

如果\emph{存在}一个 \(m\) 点集，其所有 \(2^m\) 种二值标记都能由 \(\Hcal\) 实现，就说该点集被 \(\Hcal\) shatter。用上确界而不是最大值，VC 维的严格定义为

\[
  d_{\mathrm{VC}}(\Hcal)
  =
  \sup\{m\in\mathbb{N}:\tau_{\Hcal}(m)=2^m\},
\]

Sauer 引理说明有限 VC 维会把指数级增长压成多项式增长：若 \(\operatorname{VCdim}(\Hcal)=d<\infty\)，则

\[
  \tau_{\Hcal}(m)
  \le
  \sum_{i=0}^{d}\binom{m}{i}
  \le
  \left(\frac{em}{d}\right)^d
  \quad (m\ge d).
\]

并约定集合无上界时 \(d_{\mathrm{VC}}(\Hcal)=\infty\)。这条引理的意义是：即使 \(\Hcal\) 无限，只要 VC 维有限，在 \(m\) 个样本上它能产生的标记模式也只有多项式多；但要得到随机样本上的统一收敛，不能直接把有限类证明中的 \(|\Hcal|\) 替换成 \(\tau_{\Hcal}(m)\)。

\paragraph{从 Sauer 到 VC 泛化界}

\proofsketchlabel 以下骨架针对二值类的 \(0\)--\(1\) 损失（更一般地，可替换为有界损失并相应调整常数），并取 \(S,S'\overset{\mathrm{iid}}{\sim}\Dcal^m\)。令 \(S'\) 是独立 ghost sample。先把 \(L_{\Dcal}\) 与 \(L_{S'}\) 比较，再经 symmetrization 把偏差写成 \(S\cup S'\) 上的两份经验风险之差。给定这 \(2m\) 个点后，\(\Hcal\) 只会产生至多 \(\tau_{\Hcal}(2m)\) 个标记模式，最后才对这些模式使用 union bound。因此在常数与满足 symmetrization 前提的样本量范围内，有

\[
  \Pbb\left[
    \exists h\in\Hcal:
    |L_S(h)-L_{\Dcal}(h)|>\varepsilon
  \right]
  \lesssim
  \tau_{\Hcal}(2m)\exp(-c m\varepsilon^2),
\]

其中 \(c>0\) 是集中不等式给出的常数。若 \(\operatorname{VCdim}(\Hcal)=d\)，Sauer 引理给出

\[
  \log \tau_{\Hcal}(2m)
  \le
  d\log\left(\frac{2em}{d}\right).
\]

令失败概率不超过 \(\delta\)，需要

\[
  d\log\left(\frac{2em}{d}\right)
  -
  c m\varepsilon^2
  \le
  \log\delta.
\]

移项后可读出

\[
  \varepsilon
  \gtrsim
  \sqrt{
    \frac{
      d\log(2em/d)+\log(1/\delta)
    }{m}
  }.
\]

这就是 VC 泛化界的来源。考试若要求“说明样本复杂度和 VC 维的关系”，应写出 ghost sample、symmetrization、\(\tau_{\Hcal}(2m)\) 与 Sauer 引理这条链，而不是把增长函数直接代入固定假设的 Hoeffding 界。

典型的 VC 泛化界写作

\[
  L_{\Dcal}(h)
  \le
  L_S(h)
  +
  O\left(
    \sqrt{\frac{d\log(m/d)+\log(1/\delta)}{m}}
  \right)
\]

对所有 \(h\in\Hcal\) 同时成立。考试中不必执着于常数，但必须看懂 \(d/m\) 是主导量：样本量远大于 VC 维时，训练误差才会可靠。

\subsection{Rademacher complexity}

Rademacher complexity 是另一种更细的复杂度度量。它问的是：函数类能多好地拟合随机噪声？给定样本 \(S=(z_1,\ldots,z_m)\)，经验 Rademacher complexity 为

\[
  \widehat{\mathfrak{R}}_S(\mathcal{G})
  =
  \E_{\sigma}
  \left[
    \sup_{g\in\mathcal{G}}
    \frac1m
    \sum_{i=1}^{m}\sigma_i g(z_i)
  \right],
\]

其中 \(\sigma_i\) 独立均匀取值于 \(\{-1,+1\}\)。如果 \(\mathcal{G}\) 能让内积 \(\sum_i\sigma_i g(z_i)\) 很大，说明它可以跟随机标签高度相关，复杂度高，泛化风险也高。它比 VC 维更数据依赖，在 margin 类和核方法中尤其自然。

\paragraph{Rademacher 界的推导入口}

设 \(\mathcal{G}\) 是取值于 \([0,1]\) 的损失函数类。我们想控制

\[
  \sup_{g\in\mathcal{G}}
  \left(
    \E g(z)-\frac1m\sum_{i=1}^{m}g(z_i)
  \right).
\]

引入一个与 \(S\) 独立的 ghost sample \(S'=(z_1',\ldots,z_m')\)。因为 \(\E g(z)=\E_{S'}m^{-1}\sum_i g(z_i')\)，先把总体期望替换为 ghost sample 平均：

\[
  \E_S
  \sup_g
  \left(
    \E g-\widehat{\E}_S g
  \right)
  \le
  \E_{S,S'}
  \sup_g
  \frac1m
  \sum_{i=1}^{m}
  \left(g(z_i')-g(z_i)\right).
\]

再用 Rademacher 符号随机交换 \(z_i\) 与 \(z_i'\)：

\[
  \E_{S,S'}
  \sup_g
  \frac1m
  \sum_i
  \left(g(z_i')-g(z_i)\right)
  =
  \E_{S,S',\sigma}
  \sup_g
  \frac1m
  \sum_i
  \sigma_i
  \left(g(z_i')-g(z_i)\right).
\]

由 supremum 的次可加性，右边不超过两个相同的 Rademacher 项之和，因此得到

\[
  \E_S
  \sup_g
  \left(
    \E g-\widehat{\E}_S g
  \right)
  \le
  2\mathfrak{R}_m(\mathcal{G}).
\]

这里 \(\mathfrak{R}_m(\mathcal{G})=\E_S[\widehat{\mathfrak{R}}_S(\mathcal{G})]\) 是经验 Rademacher complexity 对样本抽样再取期望后的总体版本，\(\widehat{\E}_S g=m^{-1}\sum_i g(z_i)\) 是样本平均。

高概率版本再配合 McDiarmid 或 Hoeffding 型浓缩不等式。这一推导叫 symmetrization，是 Rademacher complexity 真正进入泛化界的地方 \parencite{mohri2018foundations}。

\paragraph{实现映射}

在代码里，\(S\) 是数据矩阵和标签向量，\(h\) 是模型对象，\(L_S\) 是训练循环返回的平均 loss，\(L_{\Dcal}\) 只能用验证集或测试集近似。理论证明里常说“独立同分布抽样”；工程上对应数据切分、随机种子、去重、避免训练测试泄漏。

\paragraph{考核内容}

\begin{itemize}
  \item 写出 PAC 或 agnostic PAC 的精确定义，区分 \(\varepsilon\)、\(\delta\) 和样本复杂度。
  \item 从 Bernoulli 指示变量和 union bound 推有限类样本复杂度。
  \item 定义 shattering、增长函数、VC 维，并用 Sauer 引理得到统一收敛形式。
  \item 解释 approximation error、estimation error 与模型选择的关系。
  \item 读懂 Rademacher complexity 的定义，并说明它为什么度量“拟合随机噪声”的能力。
\end{itemize}

\section{模型选择、偏差方差与经典监督模型}

\source{本节对应考纲“模型选择、偏差方差、线性回归、逻辑回归、决策树、最近邻、SVM、核方法、压缩感知”。基础模型参考 \textcite{shalev2014uml}；核方法参考 \textcite{mohri2018foundations}；压缩感知参考 \textcite{candes2006compressive}。}

\subsection{模型选择与偏差方差}

\concepttarget{concept:bias-variance}

模型选择的核心不是选训练误差最低的模型，而是在表达能力和泛化可靠性之间平衡。把误差分解成两类更清楚：

\[
  L_{\Dcal}(h_S)-L_{\Dcal}(h^*_{\mathrm{Bayes}})
  =
  \underbrace{
    L_{\Dcal}(h^*_{\Hcal})-L_{\Dcal}(h^*_{\mathrm{Bayes}})
  }_{\text{approximation error}}
  +
  \underbrace{
    L_{\Dcal}(h_S)-L_{\Dcal}(h^*_{\Hcal})
  }_{\text{estimation error}},
\]

其中 \(h^*_{\Hcal}\) 是 \(\Hcal\) 内真实风险最小的预测器。模型类越大，approximation error 往往越小，但 estimation error 往往越大。模型类越小，估计更稳定，但可能欠拟合。

偏差方差分解是平方损失下的对应版本。设训练集随机，学习算法输出 \(\hat f_S(x)\)，对固定 \(x\) 有

\[
  \E_S[(\hat f_S(x)-f(x))^2]
  =
  \left(\E_S[\hat f_S(x)]-f(x)\right)^2
  +
  \E_S\left[
    \left(\hat f_S(x)-\E_S[\hat f_S(x)]\right)^2
  \right].
\]

第一项是 bias squared，表示平均预测离真函数多远；第二项是 variance，表示换一批训练集后预测波动多大。正则化、交叉验证、结构风险最小化（SRM）都可以看作在这两者之间调参。

\paragraph{偏差方差分解的推导}

令

\[
  \bar f(x)=\E_S[\hat f_S(x)].
\]

对误差项加减 \(\bar f(x)\)：

\[
  \hat f_S(x)-f(x)
  =
  \left(\hat f_S(x)-\bar f(x)\right)
  +
  \left(\bar f(x)-f(x)\right).
\]

平方并对训练集随机性取期望：

\[
  \E_S[(\hat f_S(x)-f(x))^2]
  =
  \E_S[(\hat f_S(x)-\bar f(x))^2]
  +
  (\bar f(x)-f(x))^2
  +
  2(\bar f(x)-f(x))
  \E_S[\hat f_S(x)-\bar f(x)].
\]

最后一项为零，因为

\[
  \E_S[\hat f_S(x)-\bar f(x)]
  =
  \E_S[\hat f_S(x)]-\bar f(x)
  =
  0.
\]

所以得到 bias squared 加 variance。若观测标签还含噪声 \(y=f(x)+\xi\)，且 \(\E[\xi\mid x]=0\)、\(\Var(\xi\mid x)=\sigma^2\)，预测 \(y\) 的均方误差还会多出不可约噪声 \(\sigma^2\)。

\subsection{线性回归}

线性回归假设

\[
  h_w(x)=\langle w,x\rangle+b,
\]

常用平方损失

\[
  \ell(h_w,(x,y))=(h_w(x)-y)^2.
\]

把样本写成设计矩阵 \(X\in\R^{m\times d}\)，其中第 \(i\) 行是 \(x_i^\top\)；标签向量写成 \(\mathbf{y}=(y_1,\ldots,y_m)^\top\in\R^m\)。忽略偏置或把偏置并入一列常数特征，最小二乘问题为

\[
  \min_{w\in\R^d}
  \frac1m\|Xw-\mathbf{y}\|_2^2.
\]

一阶条件给出正规方程

\[
  X^\top Xw=X^\top \mathbf{y}.
\]

若 \(X^\top X\) 可逆，则

\[
  w=(X^\top X)^{-1}X^\top \mathbf{y}.
\]

若不可逆，说明特征共线或样本不足，需要伪逆或正则化。岭回归把目标改成

\[
  \min_w
  \frac1m\|Xw-\mathbf{y}\|_2^2+\lambda\|w\|_2^2,
\]

一阶条件变成

\[
  (X^\top X+m\lambda I)w=X^\top \mathbf{y}.
\]

这解释了正则化的数值意义：它把病态矩阵往正定方向推，提高稳定性。
这里 \(I\in\R^{d\times d}\) 是单位矩阵，\(\lambda\ge0\) 是正则化强度。

\paragraph{正规方程的推导}

设

\[
  F(w)=\frac1m\|Xw-\mathbf{y}\|_2^2
  =
  \frac1m(Xw-\mathbf{y})^\top(Xw-\mathbf{y}).
\]

展开得到

\[
  F(w)
  =
  \frac1m
  \left(
    w^\top X^\top Xw
    -
    2\mathbf{y}^\top Xw
    +
    \mathbf{y}^\top\mathbf{y}
  \right).
\]

对 \(w\) 求梯度：

\[
  \nabla F(w)
  =
  \frac{2}{m}X^\top Xw
  -
  \frac{2}{m}X^\top\mathbf{y}.
\]

令 \(\nabla F(w)=0\)，即得到 \(X^\top Xw=X^\top\mathbf{y}\)。岭回归只是在梯度中多出 \(2\lambda w\)，所以

\[
  \frac{2}{m}X^\top Xw
  -
  \frac{2}{m}X^\top\mathbf{y}
  +
  2\lambda w
  =
  0,
\]

等价于

\[
  (X^\top X+m\lambda I)w=X^\top\mathbf{y}.
\]

\subsection{逻辑回归}

逻辑回归用于二分类，但输出的是概率：

\[
  p_w(y=1\mid x)=\sigma(\langle w,x\rangle+b),
  \qquad
  \sigma(t)=\frac{1}{1+\exp(-t)}.
\]

若标签 \(y\in\{-1,+1\}\)，逻辑损失为

\[
  \ell(w,(x,y))
  =
  \log\left(1+\exp(-y\langle w,x\rangle)\right).
\]

这可以从极大似然推出。它的优点是 convex 且可微，因此梯度下降能稳定求解。梯度为

\[
  \nabla_w\ell(w,(x,y))
  =
  -\frac{y x}{1+\exp(y\langle w,x\rangle)}.
\]

考试常把逻辑回归和 SVM 比较：逻辑损失平滑、概率解释强；hinge loss 更直接追求 margin；二者都是凸 surrogate loss，用可优化的损失替代 \(0\)-\(1\) 损失。

\paragraph{逻辑损失的 MLE 推导}

先用 \(y\in\{0,1\}\) 表示标签。设

\[
  p_w(y=1\mid x)=\sigma(\langle w,x\rangle+b),
  \qquad
  p_w(y=0\mid x)=1-\sigma(\langle w,x\rangle+b).
\]

单样本似然可以合并写成

\[
  p_w(y\mid x)
  =
  p^y(1-p)^{1-y},
  \qquad
  p=\sigma(\langle w,x\rangle+b).
\]

负对数似然为

\[
  -\log p_w(y\mid x)
  =
  -y\log p-(1-y)\log(1-p).
\]

若改用 \(y\in\{-1,+1\}\)，令 \(s=y(\langle w,x\rangle+b)\)。当 \(y=1\) 时，负对数似然是

\[
  -\log\sigma(\langle w,x\rangle+b)
  =
  \log(1+\exp(-s)).
\]

当 \(y=-1\) 时，正确类别概率为

\[
  1-\sigma(\langle w,x\rangle+b)
  =
  \sigma(-\langle w,x\rangle-b),
\]

负对数似然同样是 \(\log(1+\exp(-s))\)。因此统一得到

\[
  \ell(w,(x,y))
  =
  \log(1+\exp(-y(\langle w,x\rangle+b))).
\]

再令 \(s=y\langle w,x\rangle\) 并暂时忽略 \(b\)，链式法则给出

\[
  \nabla_w \ell
  =
  \frac{1}{1+\exp(-s)}
  \exp(-s)(-y x)
  =
  -\frac{y x}{1+\exp(s)}.
\]

这就是前面梯度公式。

\subsection{决策树}

决策树把输入空间递归切分成若干区域，每个叶子给一个预测。它的优势是可解释，缺点是若不限制深度或叶子数，容易过拟合。一个有 \(k\) 个叶子的树大致能记住 \(k\) 个局部区域，因此复杂度随树大小增长。

常见分裂准则是信息增益或 Gini impurity。若节点中正类比例为 \(p\)，熵为

\[
  H(p)=-p\log p-(1-p)\log(1-p),
\]

Gini impurity 为

\[
  G(p)=2p(1-p).
\]

分裂的目标是让子节点更纯。对候选分裂 \(a\)，信息增益可写成

\[
  \operatorname{Gain}(a)
  =
  H(S)
  -
  \sum_{v}
  \frac{|S_v|}{|S|}
  H(S_v).
\]

树模型考核通常不要求手推复杂证明，而是要求你说明过拟合来源、剪枝或 MDL/正则化思想，以及为什么贪心分裂不保证全局最优。

\subsection{最近邻算法}

\(k\)-nearest neighbors 不显式训练参数，而是在预测时寻找距离最近的训练样本。分类规则为

\[
  \hat y(x)
  =
  \operatorname{majority}
  \{y_i:x_i\in N_k(x)\},
\]

其中 \(N_k(x)\) 是距离 \(x\) 最近的 \(k\) 个训练点。\(k\) 小时方差大，容易被噪声影响；\(k\) 大时偏差大，局部结构被平均掉。距离度量和特征缩放极其重要，因为欧氏距离会被量纲大的特征支配。

\subsection{SVM、margin 与核方法}

\concepttarget{concept:svm-kernel}

线性 SVM 的核心是 margin。给定二分类样本 \(y_i\in\{-1,+1\}\)，硬间隔 SVM 求

\[
  \min_{w,b}\frac12\|w\|_2^2
  \quad
  \text{s.t.}
  \quad
  y_i(\langle w,x_i\rangle+b)\ge 1,\quad i=1,\ldots,m.
\]

为什么最小化 \(\|w\|\) 等价于最大化 margin？因为点到超平面的距离为

\[
  \frac{|\langle w,x\rangle+b|}{\|w\|_2}.
\]

约束把函数间隔固定为至少 \(1\)，于是几何间隔至少为 \(1/\|w\|_2\)。最小化 \(\|w\|\) 就是最大化间隔。

非可分数据用软间隔或 hinge loss：

\[
  \min_{w,b}
  \lambda\|w\|_2^2
  +
  \frac1m
  \sum_{i=1}^{m}
  \max\{0,1-y_i(\langle w,x_i\rangle+b)\}.
\]

核方法把输入映射到高维特征空间 \(\Phi(x)\)，但不显式计算 \(\Phi\)，只计算

\[
  K(x,x')=\langle \Phi(x),\Phi(x')\rangle.
\]

严格地说，实值核 \(k:\mathcal{X}\times\mathcal{X}\to\R\) 必须对称，即 \(k(x,x')=k(x',x)\)，并对任意有限点集 \(x_1,\ldots,x_m\) 以及任意系数 \(c\in\R^m\) 满足

\[
  \sum_{i,j=1}^{m}c_i c_j k(x_i,x_j)\ge 0.
\]

等价地，每一个有限 Gram 矩阵

\[
  [K(x_i,x_j)]_{i,j=1}^{m}
\]

半正定。在这一条件下，存在某个 Hilbert 空间及特征映射 \(\Phi\)，使 \(k(x,x')=\langle\Phi(x),\Phi(x')\rangle\)。这不是“任意相似度函数都可用”的许可；Mercer 型连续性与紧性假设只在需要积分算子谱展开时才额外出现。核技巧的考试重点是两点：第一，非线性边界可以变成高维空间线性分隔；第二，泛化不单由高维维数决定，margin 与正则化才是关键。

\paragraph{SVM 对偶与 kernel trick 推导}

硬间隔 SVM 的拉格朗日函数为

\[
  \mathcal{L}(w,b,\alpha)
  =
  \frac12\|w\|_2^2
  -
  \sum_{i=1}^{m}
  \alpha_i
  \left(
    y_i(\langle w,x_i\rangle+b)-1
  \right),
\]

其中 \(\alpha_i\ge0\) 是第 \(i\) 个约束的拉格朗日乘子。对 \(w\) 和 \(b\) 求驻点条件：

\[
  \nabla_w\mathcal{L}
  =
  w-\sum_{i=1}^{m}\alpha_i y_i x_i
  =
  0,
  \qquad
  \frac{\partial\mathcal{L}}{\partial b}
  =
  -\sum_{i=1}^{m}\alpha_i y_i
  =
  0.
\]

因此

\[
  w=\sum_{i=1}^{m}\alpha_i y_i x_i,
  \qquad
  \sum_{i=1}^{m}\alpha_i y_i=0.
\]

把 \(w\) 代回拉格朗日函数，得到对偶问题

\[
  \max_{\alpha\in\R^m}
  \sum_{i=1}^{m}\alpha_i
  -
  \frac12
  \sum_{i=1}^{m}
  \sum_{j=1}^{m}
  \alpha_i\alpha_j y_i y_j
  \langle x_i,x_j\rangle
\]

约束为

\[
  \alpha_i\ge0,
  \qquad
  \sum_{i=1}^{m}\alpha_i y_i=0.
\]

注意训练目标只通过内积 \(\langle x_i,x_j\rangle\) 依赖样本。若把样本替换为特征映射 \(\Phi(x_i)\)，内积变成

\[
  \langle \Phi(x_i),\Phi(x_j)\rangle=K(x_i,x_j).
\]

于是无需显式计算高维 \(\Phi(x)\)，只要在对偶目标中把 \(\langle x_i,x_j\rangle\) 替换为 \(K(x_i,x_j)\)。这就是 kernel trick。预测时

\[
  h(x)
  =
  \operatorname{sign}
  \left(
    \sum_{i=1}^{m}
    \alpha_i y_i K(x_i,x)
    +
    b
  \right).
\]

只有 \(\alpha_i>0\) 的训练点会出现在预测式中，它们就是 support vectors。

\subsection{压缩感知}

\concepttarget{concept:compressed-sensing}

压缩感知回答一个看似不可能的问题：如果未知信号 \(x\in\R^N\) 维度很高，但只有 \(s\) 个非零或近似非零分量，能否用远少于 \(N\) 个线性测量恢复它？

测量模型为

\[
  y=Mx,
  \qquad
  M\in\R^{K\times N},
  \qquad
  K\ll N.
\]

直接求最稀疏解是

\[
  \min_z \|z\|_0
  \quad
  \text{s.t.}
  \quad
  Mz=y,
\]

但这是组合优化。压缩感知的核心结果是，在随机测量或满足 null-space property，或满足具体定理所要求的足够强 RIP 条件的测量矩阵下，可以用凸松弛

\[
  \min_z \|z\|_1
  \quad
  \text{s.t.}
  \quad
  Mz=y
\]

精确或稳定恢复稀疏信号。RIP 的形式为：对所有 \(s\)-稀疏向量 \(z\)，

\[
  (1-\delta_S)\|z\|_2^2
  \le
  \|Mz\|_2^2
  \le
  (1+\delta_S)\|z\|_2^2.
\]

这表示 \(M\) 在所有稀疏子空间上近似保持距离，因此不会把两个不同的稀疏向量压到同一个测量。考纲中把压缩感知放在机器学习理论内，通常是考“稀疏性 + 凸松弛 + 随机测量 + 样本/测量复杂度”的基本思想，而不是完整调和分析。

\paragraph{RIP 为什么给唯一性}

设 \(z_1,z_2\) 都是 \(s\)-稀疏向量，且有同样测量：

\[
  Mz_1=Mz_2.
\]

令 \(u=z_1-z_2\)。两个 \(s\)-稀疏向量之差至多是 \(2s\)-稀疏，因此若 \(M\) 满足 \(2s\)-RIP，且 \(\delta_{2s}<1\)，则

\[
  (1-\delta_{2s})\|u\|_2^2
  \le
  \|Mu\|_2^2.
\]

但 \(Mu=Mz_1-Mz_2=0\)，于是

\[
  (1-\delta_{2s})\|u\|_2^2\le 0.
\]

由于 \(1-\delta_{2s}>0\)，只能有 \(u=0\)，即 \(z_1=z_2\)。所以 \(2s\)-RIP 至少保证了 \(\ell_0\) 问题中 \(s\)-稀疏解的可识别性；它本身并不自动推出 \(\ell_1\) 解的恢复。后者通常由 null-space property，或足够强的 RIP 条件推出。带噪观测 \(y=Mx+e\)、\(\|e\|_2\le\eta\) 时，BPDN

\[
  \min_z\|z\|_1\quad\text{s.t.}\quad\|Mz-y\|_2\le\eta
\]

在相应 RIP 条件下满足 \(\|\hat x-x\|_2\le C_0\eta+C_1\sigma_s(x)_1/\sqrt{s}\)，其中 \(\sigma_s(x)_1=\inf_{\|z\|_0\le s}\|x-z\|_1\)。这才是“稳定且近似稀疏”的精确定义 \parencite{candes2006compressive}。

\paragraph{考核内容}

\begin{itemize}
  \item 用 approximation error 和 estimation error 解释模型选择。
  \item 推线性回归正规方程，并说明不可逆时为什么需要正则化。
  \item 推逻辑回归损失和梯度，比较 logistic loss 与 hinge loss。
  \item 解释决策树为何过拟合，以及剪枝、深度限制、MDL 的作用。
  \item 写出 SVM hard/soft margin 目标，解释 margin 与 \(\|w\|\) 的关系。
  \item 定义核函数和 Gram 矩阵半正定条件。
  \item 对压缩感知说明 \(K\ll N\)、稀疏性、\(\ell_1\) 松弛和 RIP 的关系。
\end{itemize}



\section{自测与复习提纲}

\section{第三章自测：定义、条件与反例}

以下问题用于检查是否真正掌握了“对象--假设--结论”链条。每题应先写定义，再列出所需条件，最后给出一行反例或失败模式；这样写出的答案也更接近资格考试的评分点。

\begin{enumerate}
  \item \textbf{泛化与模型选择。} 为什么有限类的 union bound 不能直接替代 VC 证明？请写出 ghost sample 的作用；再区分类内最优风险、近似误差和近似 ERM 的优化误差。
  \item \textbf{经典模型。} 在线性回归中何时正规方程不可逆？岭回归怎样改变该问题？给出一个在高维未缩放特征下 \(k\)-NN 距离失真的例子。
  \item \textbf{核与压缩感知。} 为什么“某个 Gram 矩阵半正定”不足以定义核？\(2s\)-RIP 保证的唯一性是什么，为什么它仍不足以单独证明 \(\ell_1\) 恢复？
  \item \textbf{优化。} 强凸性与极小点存在性分别提供什么？请比较 GD、SGD 和随机坐标下降的随机来源，并列出 Adam 的一阶矩、二阶矩、偏差修正与数值稳定项。
  \item \textbf{深度学习。} 万能逼近定理没有保证哪两件事？任选视觉任务，写出其输入、输出、训练损失与一个常见失效模式。
  \item \textbf{强化学习。} Bellman 最优算子在哪个范数下是压缩映射？表格型 Q-learning 的收敛需要哪些访问、步长与噪声条件？
  \item \textbf{NLP/LLM。} 写出 attention logits、causal mask 与 padding mask 的区别；再说明预训练、SFT、偏好优化分别改变什么目标，以及 RAG 不能自动消除的失败模式。
\end{enumerate}

\section{从第三章到真题的使用方法}

第三章不是单独背诵材料，而是后续真题解答的工具箱。遇到题目时，先判断它属于哪门核心课，再调用对应的定义和推导模板。

\begin{longtable}{L{0.24\textwidth}L{0.34\textwidth}L{0.34\textwidth}}
\toprule
考纲主题 & 必备工具 & 常见考法 \\
\midrule
\endfirsthead
\toprule
考纲主题 & 必备工具 & 常见考法 \\
\midrule
\endhead
PAC/VC & Bernoulli 指示变量、Hoeffding、union bound、Sauer 引理 & 推样本复杂度，证明统一收敛，计算或界定 VC 维 \\
模型选择 & approximation/estimation error、bias-variance、validation & 解释欠拟合/过拟合，设计正则化或交叉验证方案 \\
线性/逻辑模型 & 正规方程、MLE、logistic loss、梯度 & 推目标函数、梯度、凸性和正则化效果 \\
SVM/核 & margin、hinge loss、KKT 直觉、PDS kernel & 写 hard/soft margin，解释核技巧和 margin 泛化 \\
压缩感知 & 稀疏性、欠定测量、\(\ell_1\) 松弛、RIP & 解释为什么少量随机测量可恢复稀疏信号 \\
凸优化 & 一阶条件、smoothness、强凸性、次梯度、prox & 证明下降、收敛率，推 prox 或 SGD 递推 \\
深度学习 & 反向传播、归一化、残差、生成模型目标 & 推梯度，比较 VAE/GAN/flow/diffusion，解释训练稳定性 \\
强化学习 & MDP、Bellman、TD error、policy gradient & 推 Bellman 方程、Q-learning、baseline 无偏性 \\
NLP/LLM & n-gram、embedding、attention、pretraining、alignment、RAG & 写 attention 维度，解释 SFT/RLHF/DPO、RAG 失败模式 \\
\bottomrule
\end{longtable}

\begin{checkpoint}
第三章的复习顺序建议是：先掌握风险与泛化，再掌握凸优化；然后学经典模型和深度网络；最后学 RL 与 NLP。原因是后两者的大多数公式都在复用前两部分的概率、优化和函数逼近语言。
\end{checkpoint}

\examnote
做真题时可以按以下模板组织答案：

\begin{enumerate}
  \item 写清对象：空间、随机变量、参数、损失或值函数。
  \item 写目标：经验风险、期望回报、似然、ELBO、Bellman residual 或对齐目标。
  \item 写关键性质：凸性、无偏性、集中界、递归关系、可逆性、mask 结构。
  \item 写结论：样本复杂度、梯度、更新式、泛化解释或模型局限。
  \item 若是开放题，补工程风险：数据泄漏、分布漂移、检索失败、优化不稳定、计算成本。
\end{enumerate}
